\documentclass[11pt]{article}

\usepackage[utf8]{inputenc}
\usepackage[T1]{fontenc}
\usepackage[margin=1in]{geometry}
\usepackage{amsmath, amssymb, amsthm}
\usepackage{microtype}
\usepackage[hidelinks]{hyperref}

\theoremstyle{plain}
\newtheorem{theorem}{Theorem}[section]
\newtheorem{lemma}[theorem]{Lemma}
\newtheorem{proposition}[theorem]{Proposition}
\newtheorem{corollary}[theorem]{Corollary}
\theoremstyle{definition}
\newtheorem{definition}[theorem]{Definition}
\theoremstyle{remark}
\newtheorem*{remark}{Remark}

\newcommand{\R}{\mathbb{R}}
\newcommand{\Q}{\mathbb{Q}}

\title{Heron's method for $\sqrt{a}$:\\[2pt]
       \large a machine-checked proof, transported to paper}
\author{R\'eflexions seminar\\ \small IECL, AHP \& Loria --- Universit\'e de Lorraine}
\date{\today}

\begin{document}
\maketitle

\begin{abstract}
We compute $\sqrt{a}$ by the Heron (Babylonian) iteration $x_{n+1}=(x_n+a/x_n)/2$ and
prove that it converges, with an explicit geometric rate and, near the root,
quadratically. The development is a machine-checked Lean~4 / Mathlib proof; this note is
its human-readable rendering, illustrating how a formal argument transports to an ordinary
mathematical text. Since the iteration uses only the field operations, the same recursion
runs as an \emph{exact, certified} square-root calculator over $\Q$.
\end{abstract}

\section{Introduction}

Fix a real number $a>0$. Heron's method improves a guess $x$ for $\sqrt{a}$ by averaging it
with $a/x$: if $x$ overestimates $\sqrt{a}$ then $a/x$ underestimates it, and their mean is
a better guess. We make this precise below, proving a geometric error bound
(Theorem~\ref{thm:geom}), quadratic convergence near the root (Proposition~\ref{prop:quad}),
and the limit itself (Theorem~\ref{thm:lim}).

Every statement corresponds to a declaration in a Lean~4 development checked against
Mathlib, whose name we give in parentheses. The point of the note is illustrative: the
formal proof and the text below carry the \emph{same} content, and the passage from one to
the other is essentially mechanical --- proof-assistant plumbing disappears, and the
load-bearing identities become the displayed equations.

\section{The iteration}

\begin{definition}[Heron step and sequence; \texttt{step}, \texttt{heron}]
For $x\neq 0$ set
\[
  S(x)\;=\;\frac{x+a/x}{2}.
\]
Given $x_0>0$, define $x_0,x_1,x_2,\dots$ by $x_{n+1}=S(x_n)$.
\end{definition}

\begin{remark}
$S$ is \emph{computable}, over any field --- in particular over $\Q$, which we use in Section~\ref{sec:calc}.
The proofs themselves take place over $\R$.
\end{remark}

\begin{lemma}[Rational form; \texttt{step\_eq}]\label{lem:ratform}
For $x\neq 0$,\qquad $\displaystyle S(x)=\frac{x^2+a}{2x}.$
\end{lemma}

\begin{proof}
Multiplying through by $x$, $\dfrac{x+a/x}{2}=\dfrac{x^2+a}{2x}$.
\end{proof}

\begin{lemma}[Positivity; \texttt{step\_pos}, \texttt{heron\_pos}]\label{lem:pos}
If $x>0$ then $S(x)>0$; consequently $x_n>0$ for every $n$.
\end{lemma}

\begin{proof}
If $x>0$ then $a/x>0$, so $S(x)=(x+a/x)/2>0$. Since $x_0>0$, induction gives $x_n>0$ for
all $n$.
\end{proof}

\section{Convergence}

\begin{lemma}[The step overshoots; \texttt{le\_sq\_step}]\label{lem:overshoot}
For $x\neq 0$,\qquad $a\leq S(x)^2.$
\end{lemma}

\begin{proof}
By Lemma~\ref{lem:ratform},
\[
  S(x)^2-a=\Big(\frac{x^2+a}{2x}\Big)^2-a=\Big(\frac{x^2-a}{2x}\Big)^2\;\geq\;0. \qedhere
\]
\end{proof}

\begin{corollary}[Lower bound; \texttt{sqrt\_le\_step}, \texttt{sqrt\_le\_heron}]\label{cor:lower}
If $x>0$ then $\sqrt{a}\leq S(x)$. In particular $\sqrt{a}\leq x_n$ for every $n\geq 1$.
\end{corollary}

\begin{proof}
Both sides are nonnegative and $a\leq S(x)^2$ by Lemma~\ref{lem:overshoot}, so
$\sqrt{a}\leq\sqrt{S(x)^2}=S(x)$. Taking $x=x_n$ gives $\sqrt{a}\leq x_{n+1}$, that is,
$\sqrt{a}\leq x_n$ for $n\geq 1$.
\end{proof}

\begin{lemma}[Error identity; \texttt{step\_sub\_sqrt}]\label{lem:err}
For $x>0$,\qquad $\displaystyle S(x)-\sqrt{a}=\frac{(x-\sqrt{a})^2}{2x}.$
\end{lemma}

\begin{proof}
Using Lemma~\ref{lem:ratform} and $(\sqrt{a})^2=a$,
\[
  S(x)-\sqrt{a}=\frac{x^2+a-2\sqrt{a}\,x}{2x}
              =\frac{x^2-2\sqrt{a}\,x+(\sqrt{a})^2}{2x}
              =\frac{(x-\sqrt{a})^2}{2x}. \qedhere
\]
\end{proof}

\begin{lemma}[Contraction; \texttt{step\_sub\_sqrt\_le}]\label{lem:contr}
If $\sqrt{a}\leq x$, then $\displaystyle S(x)-\sqrt{a}\leq\frac{x-\sqrt{a}}{2}.$
\end{lemma}

\begin{proof}
By Lemma~\ref{lem:err}, the gap between the two sides is
\[
  \frac{x-\sqrt{a}}{2}-\big(S(x)-\sqrt{a}\big)
   =\frac{x-\sqrt{a}}{2}-\frac{(x-\sqrt{a})^2}{2x}
   =\frac{\sqrt{a}\,(x-\sqrt{a})}{2x}\;\geq\;0,
\]
since $x>0$ and $x\geq\sqrt{a}\geq 0$.
\end{proof}

\begin{theorem}[Geometric convergence; \texttt{heron\_error\_le}]\label{thm:geom}
For every $k\geq 0$,
\[
  x_{k+1}-\sqrt{a}\;\leq\;\Big(\tfrac12\Big)^k\,(x_1-\sqrt{a}).
\]
\end{theorem}

\begin{proof}
We argue by induction on $k$, the case $k=0$ being an equality. For the inductive step,
$x_{k+1}\geq\sqrt{a}$ by Corollary~\ref{cor:lower}, so the contraction
(Lemma~\ref{lem:contr}) followed by the induction hypothesis gives
\[
  x_{k+2}-\sqrt{a}\;\leq\;\frac{x_{k+1}-\sqrt{a}}{2}
   \;\leq\;\frac12\Big(\tfrac12\Big)^k(x_1-\sqrt{a})
   \;=\;\Big(\tfrac12\Big)^{k+1}(x_1-\sqrt{a}). \qedhere
\]
\end{proof}

\section{Sharper results}

\begin{proposition}[Quadratic convergence; \texttt{step\_sub\_sqrt\_le\_sq}]\label{prop:quad}
If $\sqrt{a}\leq x$, then $\displaystyle S(x)-\sqrt{a}\leq\frac{(x-\sqrt{a})^2}{2\sqrt{a}}.$
\end{proposition}

\begin{proof}
By Lemma~\ref{lem:err}, $S(x)-\sqrt{a}=(x-\sqrt{a})^2/(2x)$. Since $0<\sqrt{a}\leq x$ we
have $2\sqrt{a}\leq 2x$, hence enlarging the denominator gives
$(x-\sqrt{a})^2/(2x)\leq(x-\sqrt{a})^2/(2\sqrt{a})$.
\end{proof}

\begin{remark}
This is the source of the method's speed: once $x_n$ is close to $\sqrt{a}$, the error is
essentially squared at each step, so the number of correct digits roughly doubles per
iteration.
\end{remark}

\begin{theorem}[Limit; \texttt{heron\_tendsto}]\label{thm:lim}
$x_n\to\sqrt{a}$ as $n\to+\infty$.
\end{theorem}

\begin{proof}
For $n\geq 1$, Corollary~\ref{cor:lower} and Theorem~\ref{thm:geom} give
\[
  0\;\leq\;x_{n+1}-\sqrt{a}\;\leq\;\Big(\tfrac12\Big)^n(x_1-\sqrt{a}),
\]
and the right-hand side tends to $0$; the claim follows by squeezing.
\end{proof}

\section{A certified calculator}\label{sec:calc}

\begin{remark}[\texttt{sqrtApprox}, \texttt{le\_sq\_step\_rat}]
Because $S$ uses only the field operations, the recursion runs verbatim over $\Q$, where it
is exactly computable. The overshoot inequality $a\leq S(x)^2$ (Lemma~\ref{lem:overshoot})
holds over $\Q$ as well, so from step~$1$ on, every computed iterate is a rational whose
square exceeds $a$ --- a \emph{certified} upper bound for $\sqrt{a}$. Iterating until
$\vert x^2-a\vert\leq\varepsilon$ then returns $\sqrt{a}$ to any requested precision; for
instance four steps from $x_0=1$ give
\[
  \sqrt{2}\;\approx\;\frac{665857}{470832}\;=\;1.4142135623\ldots,
\]
already correct to ten places. In this sense the formal proof yields a certified
square-root calculator: the program that runs is exactly the function the theorems above
describe.
\end{remark}

\end{document}
